% Introduction section

% This section introduces the background and context for the study.

\subsection{Background}
Mendelian randomization (MR) is an analytical method that uses genetic variants (e.g. single-nucleotide polymorphisms (SNPs)) as instrumental variables (IVs) to infer causal relationships between exposures and outcomes, minimizing unmeasured confounding that often bias traditional observational studies \cite{smith_mendelian_2003, Nguyen2024, Davies2018}. Over the past two decades, extensions of MR methods\cite{Burgess2015, burgess_guidelines_2023, burgess_network_2015}, software packages\cite{hemani_mr-base_2018}, and public datasets\cite{sudlow_uk_2015} have emerged. And the number of published MR studies has increased substantially\cite{He2025}, reflecting the growing popularity of the approach in epidemiology and biomedical research. 

However, concerns have been raised regardsing the completeness, transparency, and reproducibility of MR reporting. Many studies provide insufficient methodological details, which limits their interpretability and undermines the credibility of causal claims \cite{davies_issues_2013,islam_reporting_2022}. For example, IVs used in a MR study must meet three core assumptions (relevance, independence, and exclusion)\cite{VanderWeele2014}. So naturally, people could ask: "Had authors assessed these assumptions in their report?" or "Had authors tested the robustness of their results against violations of assumptions?". As two-sample MR studies became more and more popular, questions can also be raised on the similarity of the genetic variant-exposure associations between the exposure and outcome samples. To address these issues, the STROBE-MR (Strengthening the Reporting of Observational Studies in Epidemiology for Mendelian Randomization) guideline was developed. The STROBE-MR guideline, which contains 43 subitems (we use "items" for simplicity in the remainder of this paper) that cover 6 sections (Title and Abstract, Introduction, Methods, Results, Discussion and Other Information), provides a structured framework to assist authors in reporting MR studies clearly and help reviewers to assess the quality of MR studies \cite{skrivankova_strengthening_2021}. 

Despite the guideline's availability, assessments of MR study quality remain largely manual, relying on expert reviewers to evaluate adherence to STROBE-MR criteria \cite{li_causal_2025}. This manual approach limits large-scale assessment of the MR reporting. Consequently, there is a need for automated methods that can efficiently and accurately evaluate reporting quality at scale. 

Recent advances in large language models (LLMs) offer new opportunities: these models can process and interpret complex biomedical texts efficiently, demonstrating strong potential in biomedical knowledge extraction\cite{legate_semiautomated_2024}, literature selection\cite{laignelot_large_2026}, risk of bias assessment\cite{nyrhi_large_2026, lai_assessing_2024}, and reporting quality assessment\cite{srinivasan_large_2025, luo_using_2025}. In this study, we developed an LLM based approach to address challenges of large-scale report quality assessment and evaluated its performance on an annotated dataset. We then applied this approach to over 13, 000 MR studies published between 2007 and 2025 to analyze trends of reporting in different items, sections, and time periods.

We first describe the collection of MR studies, strategies we used to divide STROBE-MR items into smaller components, prompt design and model evaluation process. We then present the model performance and states of MR reporting. Finnaly, we discuss implications for how we can improve the reporting quality and prompt design for assessing MR reporting or similar tasks.

\subsection{Research objectives}
This project aims to leverage LLMs to develop and validate an automated system capable of rapidly assessing reporting quality of MR publications based on STROBE-MR guideline. Such a system would enable us to explore and describe the general state of MR reporting and provide insights on better reporting practices and more reliable causal inference in epidemiology. 

% [Research questions and objectives]

% Example cross-reference to supplementary materials
% Additional background information is provided in Section~\ref{sec:supp-methods} of the supplementary materials.